% TOP_PROBLEM: {"id": 30006870, "title": "Shelah's Conjecture on NIP Fields", "category_id": 18, "category": {"id": 18, "name": "logic", "display_name": "Logic", "description": "Problems in mathematical logic, model theory, proof theory, and finite model theory.", "slug": "logic", "order_index": 18, "created_at": "2026-07-31T15:26:25.671Z"}, "status": "open"}
% FIRST_POSTED: 2026-09-25
% ENTRY_KIND: statement_only
% !TeX program = lualatex
% Definition and sources only; this file is not an LLM solution attempt.
\documentclass[11pt]{article}
\usepackage[margin=25mm]{geometry}
\usepackage{fontspec}
\setmainfont{Latin Modern Roman}
\usepackage{amsmath,amssymb,mathtools}
\usepackage{unicode-math}
\setmathfont{Latin Modern Math}
\usepackage{xurl,hyperref}
\hypersetup{colorlinks=true,urlcolor=blue}
\setcounter{secnumdepth}{0}
\setlength{\parindent}{0pt}
\setlength{\parskip}{5pt}
\setlength{\emergencystretch}{3em}
\begin{document}
\section*{\texorpdfstring{Shelah's Conjecture on NIP Fields}{Shelah's Conjecture on NIP Fields}}\label{30006870}
\subsection{Definitions and mathematical statement}
A formula $\varphi(x;y)$ has the independence property if, in models of the theory, it can realize every subset of arbitrarily large finite sets: there are $a_1,\ldots,a_n$ and parameters $b_S$ with $\varphi(a_i;b_S)$ exactly when $i\in S$. A theory is NIP if no formula has this property. For an infinite field $K$ in the ring language, the selected conjecture asserts that NIP implies one of: $K$ is separably closed; $K$ is real closed; or $K$ admits a nontrivial henselian valuation. A valuation is henselian when its valuation ring satisfies Hensel's lifting of simple residue-field roots; nontrivial means its valuation ring is not all of $K$. No definability of the valuation is demanded by this formulation. The all-NIP-field quantifier is not replaced by finite-rank or otherwise restricted subclasses.
\par\smallskip\textit{Formulation and concept references:} \href{https://arxiv.org/abs/1911.00309}{[S1]}.

\subsection{Short English statement}
Must every infinite NIP field be separably closed, real closed, or carry a nontrivial henselian valuation?
\subsection{Sources}
\begin{itemize}
\item[S1]Characterizing NIP henselian fields (arXiv abstract). Accessed 2026-09-01.
\url{https://arxiv.org/abs/1911.00309}.
\textit{Formulation source.}
\end{itemize}
\end{document}
